% CE520 M.S. Graduate Seminar — thesis progress, Fall 2026
% Source for local TeX Live / MiKTeX. Keep the .sty files, references.bib,
% and logos/ directory alongside main.tex.
% Compile: pdflatex main.tex; biber main; pdflatex main.tex; pdflatex main.tex
% The compiled PDF is the normal, full-screen presentation format for Beamer.

\documentclass[aspectratio=43,11pt]{beamer}
\usetheme{metu}
\usefonttheme{serif}
\usepackage{lmodern}
\usepackage{amsmath,amssymb}
\usepackage{booktabs}
\usepackage{array}
\usetikzlibrary{arrows.meta,calc,positioning,decorations.pathreplacing,shapes.geometric,fit}
\usepackage{graphicx}

% ---------------- EDIT THESE FIELDS ----------------
\newcommand{\PresenterName}{Bilge Kutay}
\newcommand{\Advisor}{Asst. Prof. Dr. Cüneyt Baykal}
\newcommand{\TalkDate}{November 11, 2026}
\newcommand{\CourseCode}{CE520 Graduate Seminar}
% ---------------------------------------------------

\title{Machine-Learning-Based Active Wave Absorption in a Numerical Wave Flume}
\author{\PresenterName\\[1mm] \scriptsize M.S. Candidate, Coastal and Harbor Engineering\\[1mm] \scriptsize Supervisor: \Advisor}
\institute{CE520 Graduate Seminar -- Fall 2026}
\date{\TalkDate}
\addbibresource{references.bib}
\AtBeginBibliography{%
  \fontsize{8}{8.8}\selectfont
}
\DeclareLanguageMapping{english}{english-apa}
\renewcommand*{\nameyeardelim}{\addcomma\space}

% Typography and diagram palette
\definecolor{water}{RGB}{210,237,249}
\definecolor{deepblue}{RGB}{19,93,166}
\definecolor{teal}{RGB}{0,120,133}
\definecolor{softgray}{RGB}{246,247,249}
\definecolor{amber}{RGB}{171,113,28}
\setbeamercolor{normal text}{fg=metublack}
\setbeamerfont{normal text}{size=\small}
\setbeamerfont{block title}{size=\small,series=\bfseries}
\setbeamerfont{block body}{size=\small}
\setbeamersize{text margin left=8mm,text margin right=8mm}
\setlength{\parskip}{3pt}
\newcommand{\subtle}[1]{{\scriptsize\color{metugray}#1}}
\newcommand{\smallsource}[1]{\par\vspace{1mm}{\tiny\color{metugray}#1}}
\newcommand{\redrule}{\color{metured}\rule{\linewidth}{0.6pt}}
\newcommand{\tagline}[1]{{\color{metured}\bfseries #1}}
\newcommand{\status}[1]{\textcolor{metured}{\textbf{#1}}}
\tikzset{
  flow/.style={draw=metugray!65,fill=softgray,rounded corners=2pt,
    align=center,inner sep=4pt,font=\scriptsize},
  redflow/.style={draw=metured!85,fill=metured!7,rounded corners=2pt,
    align=center,inner sep=4pt,font=\scriptsize},
  arr/.style={-{Latex[length=2.2mm]},line width=0.85pt,deepblue},
  arrred/.style={-{Latex[length=2.2mm]},line width=0.85pt,metured}
}

\begin{document}

% ==============================================================
% 01 — approx. 15 s: Introduce the thesis as ongoing work.
% ==============================================================
\begin{frame}[noframenumbering,plain]
  \titlepage
\end{frame}

\begin{frame}[t]{Outline}
\color{metublack}
    \tableofcontents
\end{frame}

% ==============================================================
% 02 — approx. 55 s: Observed engineering problem.
% ==============================================================
\section{Problem and Motivation}
\begin{frame}[t]{Why are reflections a problem?}
\begin{center}
\begin{tikzpicture}[x=0.92cm,y=0.92cm]
  \fill[water] (0.7,0.55) rectangle (10.8,1.55);
  \fill[metugray!70] (0.58,0.44) rectangle (11.00,0.56);
  \fill[metugray!70] (10.80,0.55) rectangle (10.96,2.50);
  \fill[metured] (0.65,0.55) rectangle (0.89,2.45);
  \draw[deepblue,line width=1.1pt,samples=90,domain=0.9:10.8]
     plot (\x,{1.55+0.13*sin(360*(\x-0.9)/1.65)});
  \draw[metured,dashed,line width=1.1pt,samples=90,domain=0.9:10.8]
     plot (\x,{1.55+0.23*sin(360*(\x-3.9)/1.65)});
  \draw[arr,line width=1.8pt] (2.4,2.10)--(6.7,2.10);
  \node[deepblue,font=\scriptsize\bfseries] at (4.55,2.35) {Incident waves};
  \draw[arrred,dashed,line width=1.5pt] (9.45,1.02)--(5.75,1.02);
  \node[metured,font=\scriptsize\bfseries] at (7.57,0.77) {Reflected waves};
  \node[metured,font=\scriptsize,anchor=south] at (1.0,2.58) {Wavemaker};
  \node[metugray,font=\scriptsize,anchor=south] at (10.30,2.63) {End boundary};
\end{tikzpicture}
\end{center}
\vspace{0.5mm}
\begin{block}{Observed problem}
  Reflected waves return to the wavemaker and may be \textcolor{metured}{re-reflected},
  contaminating the intended wave field during long experiments.
\end{block}
\vspace{1mm}
\begin{itemize}\setlength{\itemsep}{2pt}
  \item Reflections alter free-surface records and limit useful test duration.
  \item Reliable experiments require \textbf{wave generation and absorption}.
\end{itemize}
\end{frame}

% ==============================================================
% 03 — approx. 60 s: Point of departure and precise research gap.
% ==============================================================
\section{Research Gap}

\begin{frame}[t]{What is known, and what is missing?}

\begin{columns}[T,onlytextwidth]

  \begin{column}{0.30\textwidth}
    \begin{block}{Passive absorption}
      \begin{minipage}[t][2.18cm][t]{\linewidth}
        \raggedright\footnotesize
        Beaches and porous layers dissipate wave energy,
        but need space and are sensitive to wave conditions.
      \end{minipage}
    \end{block}
  \end{column}

  \hfill

  \begin{column}{0.30\textwidth}
    \begin{block}{Classical AWA}
      \begin{minipage}[t][2.18cm][t]{\linewidth}
        \raggedright\footnotesize
        Wave-gauge feedback and physics-based filters
        correct wavemaker motion to limit re-reflection.
      \end{minipage}
    \end{block}
  \end{column}

  \hfill

  \begin{column}{0.30\textwidth}
    \begin{block}{ML wave control}
      \begin{minipage}[t][2.18cm][t]{\linewidth}
        \raggedright\footnotesize
        Xie and colleagues used ANNs for piston
        and plunger wave generation and absorption
        in their own CFD solver.
      \end{minipage}
    \end{block}
  \end{column}

\end{columns}

\vspace{2.5mm}

\begin{beamercolorbox}[
  rounded=true,
  sep=1.0ex,
  wd=\linewidth
]{block body}
  {\small
  \textcolor{metured}{
    \textbf{Research gap and planned extension}
  }\\[1mm]
  Implement \textbf{ML-based active absorption}
  within OpenFOAM's moving-piston system and
  \textbf{compare its performance with conventional AWA}
  under the same wave conditions.
  }
\end{beamercolorbox}

\smallsource{
  Sch\"affer \& Jakobsen (2003);
  Altomare et al. (2017);
  Xie \& Zhao (2022);
  Xie et al. (2023).
}

\end{frame}

% ==============================================================
% 04 — approx. 40 s: Research questions.
% ==============================================================
\section{Research Questions}
\begin{frame}[t]{Thesis Objective and Research Questions}
\vspace{0.5cm}
\begin{block}{Overall objective}
  Implement and evaluate a \textbf{machine-learning-based active wave absorption}
  controller within OpenFOAM's piston-type wavemaker system.
\end{block}
\vspace{2mm}
\begin{enumerate}
    \item Can a trained ML controller be integrated into
    OpenFOAM's moving-piston system for active wave absorption?
    \item How does ML-based absorption perform
    compared with conventional active absorption
    methods?
    \item How does the controller's performance
    vary under different wave conditions?
\end{enumerate}
\end{frame}

% ==============================================================
% 05 — approx. 65 s: Flume geometry, 13 gauges, important near-paddle WG0.
% ==============================================================
\section{Numerical Model}
\begin{frame}[t]{OpenFOAM Numerical Wave-flume Setup}
\begin{center}
\begin{tikzpicture}[x=0.88cm,y=0.90cm]
  % Flume geometry is a schematic; gauge distances are listed below.
  \fill[water] (0.4,0.57) rectangle (11.3,1.42);
  \draw[deepblue,very thick,samples=110,domain=0.65:11.1]
    plot (\x,{1.42+0.105*sin(360*(\x-0.65)/1.5)});
  \fill[metugray!70] (0.40,0.42) rectangle (11.33,0.58);
  \fill[metugray!70] (11.16,0.57) rectangle (11.33,2.21);
  \fill[metured] (0.57,0.58) rectangle (0.78,2.18);
  \draw[<->,metured,line width=1pt] (0.18,1.08)--(1.18,1.08);
  \node[font=\scriptsize,metured,above] at (0.65,2.30) {Piston};
  % Near-paddle gauge; labels offset so near-zero spacing is intelligible.
  % Near-paddle gauge
  \draw[teal,line width=1.0pt] (1.22,0.63)--(1.22,1.85);
  \fill[teal] (1.22,1.41) circle (0.052);
  \node[font=\tiny\bfseries,teal,anchor=south] at (1.22,1.91) {WG0};
  \draw[teal,densely dotted] (0.83,1.61)--(1.16,1.61);
  \draw[metugray!60] (1.22,0.38)--(1.22,0.25);
  \node[font=\scriptsize] at (1.22,0.06) {$0.10$ m};
  % Four triplets at nominal x = 3, 6, 9, 12 m. Schematically spaced.
  \foreach \xc/\tag/\pos in {3.12/{WG1--3}/3,5.42/{WG4--6}/6,7.72/{WG7--9}/9,10.02/{WG10--12}/12}{
     \foreach \dx in {-0.16,0,0.16}{
       \draw[teal,line width=0.9pt] (\xc+\dx,0.62)--(\xc+\dx,1.85);
       \fill[teal] (\xc+\dx,1.43) circle (0.045);
     }
     \node[font=\tiny\bfseries,teal,anchor=south] at (\xc,1.91) {\tag};
     \draw[metugray!60] (\xc,0.38)--(\xc,0.25);
     \node[font=\scriptsize] at (\xc,0.06) {$\pos$ m};
  }
\end{tikzpicture}
\end{center}
\vspace{1mm}
\begin{columns}[T,onlytextwidth]
\begin{column}{0.49\textwidth}
\begin{itemize}
    \item \textbf{Numerical model:} OpenFOAM / olaDyMFlow
    \item \textbf{Flume geometry:} $L \approx 17.5$ m, $h = 0.50$ m
    \item \textbf{Wavemaker:} Piston-type with deforming mesh
\end{itemize}
\end{column}
\begin{column}{0.53\textwidth}
\begin{itemize}
    \item \textbf{Numerical resolution:} 374,872 cells; $\Delta t = 0.0005$ s
    \item \textbf{Mesh resolution:} $\Delta x = 4$ mm near piston; $\Delta z = 2$ mm near surface
    \item \textbf{Wave case:} $T = 1.0$ s, $H \approx 0.071$ m, stroke $= 0.04$ m
\end{itemize}
\end{column}
\end{columns}
\end{frame}

% ==============================================================
% 06 — approx. 60 s: Study design and validation benchmarks.
% ==============================================================
\section{Numerical Methodology}
\begin{frame}[t]{Numerical Methodology}
\begin{center}
\begin{tikzpicture}[node distance=7mm and 6mm]
  \node[flow,minimum width=2.30cm,minimum height=1cm] (model)
    {CFD model\\verification};

  \node[flow,right=of model,minimum width=2.30cm,minimum height=1cm] (waves)
    {Prescribed\\piston motion};

  \node[flow,right=of waves,minimum width=2.30cm,minimum height=1cm] (gauge)
    {Piston-side\\water level $\eta(t)$};

  \node[redflow,right=of gauge,minimum width=2.30cm,minimum height=1cm] (eval)
    {ML training\\dataset};
  \draw[arr] (model)--(waves);
  \draw[arr] (waves)--(gauge);
  \draw[arr] (gauge)--(eval);
\end{tikzpicture}
\end{center}
\vspace{4mm}
\textbf{Comparison strategy}
\begin{enumerate}\setlength{\itemsep}{4pt}
  \item No active absorption \hfill \subtle{reference / uncontrolled case}
  \item Conventional active absorption \hfill \subtle{classical benchmark}
  \item ML-based active absorption \hfill \subtle{planned OpenFOAM implementation}
\end{enumerate}
\vspace{3mm}
\begin{block}{Wave reflection analysis}
  Use virtual wave-gauge measurements to estimate
  the reflection coefficient $K_r=H_r/H_i$
  for each of the three cases.
\end{block}
\end{frame}

% ==============================================================
% 07 — approx. 60 s: Feedback method, terms and separation.
% ==============================================================
\section{ML-based Controller}
\begin{frame}[t]{ML-based Feedback Control}
\vspace{-2mm}
\begin{center}

\begin{tikzpicture}[x=1cm,y=1cm]

  % Target wave: separate from the feedback loop
  \node[flow,minimum width=1.9cm,
        minimum height=0.95cm] (target) at (0,1.6)
    {Target wave\\$\eta_d(t)$};

  % Upper part of the feedback loop
  \node[flow,minimum width=2.3cm,
        minimum height=0.95cm] (flume) at (3.4,1.6)
    {Moving piston\\and CFD flume};

  \node[flow,minimum width=2.2cm,
        minimum height=0.95cm] (measured) at (7,1.6)
    {Measured level\\$\eta_m(t)$};

  % Lower part of the feedback loop
  \node[flow,minimum width=2.5cm,
        minimum height=1cm] (corrected) at (7,0)
    {Corrected target\\$\eta_c=2\eta_d-\eta_m$};

  \node[redflow,minimum width=2.4cm,
        minimum height=1cm] (ann) at (3.4,0)
    {Trained ANN\\predicts piston velocity};

  % Target input
  \draw[arr] (target.east) -- (flume.west);

  % Closed feedback loop
  \draw[arr] (flume.east) -- (measured.west);
  \draw[arr] (measured.south) -- (corrected.north);
  \draw[arrred] (corrected.west) -- (ann.east);
  \draw[arrred] (ann.north) -- (flume.south);

  % Feedback arrow label
  \node[font=\tiny,metured,anchor=west]
    at ($(ann.north)!0.5!(flume.south)+(0.12cm,0)$)
    {Piston update};

  
  % Wave elevation notation — below target wave
  \node[
      anchor=north west,
      align=left,
      font=\tiny,
      text=black!65
  ] at ($(target.south west)+(-1cm,-0.55cm)$) {
      $\eta_d$: Desired elevation\\[2pt]
      $\eta_m$: Measured elevation\\[2pt]
      $\eta_c$: Corrected elevation
  };

\end{tikzpicture}

\end{center}
\[
\boxed{\,X_{\mathrm{piston}}(t)=X_{\mathrm{gen}}(t)+X_{\mathrm{abs}}(t)\,}
\]
\vspace{-4mm}

\begin{itemize}\setlength{\itemsep}{3pt}
  \item \textbf{ANN inputs:} corrected target signal,
        wave period $T$, and water depth $h$.
  \item \textbf{Output:} predicted piston velocity,
        used to update piston displacement.
  \item \textbf{Control:} limit piston drift while
        maintaining the intended waves.
\end{itemize}

\end{frame}

% ==============================================================
% 09 — approx. 55 s: Metrics and expected results (not claims).
% ==============================================================
\section{Expected Results}
\begin{frame}[t]{Evaluation and Expected Contribution}
\begin{columns}[T,onlytextwidth]
\begin{column}{0.48\textwidth}
\textbf{How performance will be measured}
\begin{itemize}\setlength{\itemsep}{2pt}
  \item Agreement with target wave elevation
  \item Reflection coefficient $K_r=H_r/H_i$
  \item Performance under different wave conditions
  \item Piston-motion limits and stability
\end{itemize}
\end{column}
\begin{column}{0.49\textwidth}
\textbf{What the thesis aims to deliver}
\begin{itemize}\setlength{\itemsep}{2pt}
  \item Reproducible CFD benchmark cases
  \item Training dataset generated in OpenFOAM
  \item ML-based absorption integrated into OpenFOAM
  \item Comparison with conventional AWA
\end{itemize}
\end{column}
\end{columns}
\vspace{1mm}
\begin{block}{Project-level target}
  The larger laboratory project targets \textbf{$K_r<0.10$} under its intended test
  conditions.
\end{block}
\vspace{1mm}
% Reflection-separation methods and analysis windows will be documented.
\end{frame}

% ==============================================================
% 10 — approx. 35 s: Practical relevance.
% ==============================================================
\section{Practical Relevance}
\begin{frame}[t]{Why does this matter in practice?}
\begin{center}
\begin{tikzpicture}[node distance=5mm and 5mm]
  \node[flow,minimum height=1.3cm,minimum width=2.55cm] (a)
    {Cleaner laboratory\\wave conditions};
  \node[flow,right=of a,minimum height=1.3cm,minimum width=2.55cm] (b)
    {Longer useful\\testing windows};
  \node[flow,right=of b,minimum height=1.3cm,minimum width=2.55cm] (c)
    {More reliable\\coastal model tests};
  \draw[arr] (a)--(b);
  \draw[arr] (b)--(c);
\end{tikzpicture}
\end{center}
\vspace{2mm}
\begin{itemize}\setlength{\itemsep}{8pt}
  \item \textbf{Scientific value:}
        improve the accuracy of wave experiments by
        reducing unwanted re-reflections.
  \item \textbf{Engineering value:}
        enable longer, more reliable tests of coastal
        and offshore structures.
  \item \textbf{Open-source value:}
        demonstrate ML-based active absorption within OpenFOAM's moving-piston system.
  \item \textbf{Institutional value:}
        support the development of a wave-control
        system for METU's laboratory.
\end{itemize}
\end{frame}

% ==============================================================
% 11 — approx. 50 s: Requested Gantt chart through June 2027.
% ==============================================================
\section{Research Schedule}
\begin{frame}[t]{Research Timeline}
\vfill
\centering
\begin{tikzpicture}[x=0.95cm,y=0.80cm,font=\tiny]
  \def\xstart{3.25}
  \def\mw{0.98}

  % Month columns
  \foreach \i/\m in {0/Nov,1/Dec,2/Jan,3/Feb,4/Mar,5/Apr,6/May,7/Jun}{
    \fill[softgray] (\xstart+\i*\mw,7.45) rectangle ++(\mw,0.60);
    \draw[metugray!35] (\xstart+\i*\mw,-0.13)--++(0,8.18);
    \node[font=\scriptsize\bfseries] at (\xstart+\i*\mw+0.49,7.72) {\m};
  }
  \draw[metugray!40] (\xstart+8*\mw,-0.13)--++(0,8.18);

  % Activities
  \foreach \yy/\lab in {
    6.95/{CFD / piston verification},
    6.10/{Wave-gauge validation},
    5.25/{Wave-generation validation},
    4.40/{Baseline AWA comparison},
    3.55/{Training data preparation},
    2.70/{ML training / integration},
    1.85/{Controller evaluation},
    1.00/{Performance comparison},
    0.15/{Thesis writing / defence}
  }{
    \node[anchor=east,align=right] at (\xstart-0.15,\yy) {\lab};
    \draw[metugray!25] (-0.35,\yy-0.28)--(\xstart+8*\mw,\yy-0.28);
  }

  % Schedule bars
  \fill[deepblue!85,rounded corners=1pt] (\xstart+0.06,6.73) rectangle (\xstart+1.90,7.16);
  \fill[deepblue!85,rounded corners=1pt] (\xstart+0.06,5.88) rectangle (\xstart+1.90,6.31);
  \fill[deepblue!85,rounded corners=1pt] (\xstart+1.04,5.03) rectangle (\xstart+2.88,5.46);
  \fill[deepblue!85,rounded corners=1pt] (\xstart+1.04,4.18) rectangle (\xstart+2.88,4.61);
  \fill[teal!85,rounded corners=1pt] (\xstart+3.00,3.33) rectangle (\xstart+4.84,3.76);
  \fill[teal!85,rounded corners=1pt] (\xstart+3.98,2.48) rectangle (\xstart+5.82,2.91);
  \fill[metured!85,rounded corners=1pt] (\xstart+4.96,1.63) rectangle (\xstart+6.80,2.06);
  \fill[metured!85,rounded corners=1pt] (\xstart+5.94,0.78) rectangle (\xstart+6.80,1.21);
  \fill[metured!85,rounded corners=1pt] (\xstart+4.96,-0.08) rectangle (\xstart+7.78,0.35);
\end{tikzpicture}
\vfill
\end{frame}


% ==============================================================
% 12 — approx. 40 s: Required elevator pitch and conclusion.
% ==============================================================
\section{Conclusion}
\begin{frame}[t]{Conclusion: the thesis in three sentences}
\vspace{0.4cm}
\begin{beamercolorbox}[rounded=true,sep=2.5ex,wd=\linewidth]{block body}
  \large
  Unwanted re-reflections reduce the reliability of laboratory wave testing.

  \vspace{3mm}
  This thesis investigates whether a \textcolor{metured}{machine-learning based
  feedback controller} can adjust a piston wavemaker to absorb returning waves while maintaining accurate wave generation.

  \vspace{3mm} 
  By integrating the trained controller directly into
  \textbf{OpenFOAM} and comparing it with conventional
  active absorption, this study aims to evaluate its
  feasibility and performance.
\end{beamercolorbox}
\vfill
\end{frame}

% ==============================================================
% 13 — Backup references. Does not need discussion in the 8-10 min.
% ==============================================================
\begin{frame}[t]{Selected References}
\nocite{schaffer2003,altomare2017,xie2023,xie2022easily,higuera2013realistic,schaffer2000review,gillan2022}
\printbibliography[heading=none]
\end{frame}

% ==============================================================
% 14 — final Q&A.
% ==============================================================
\makethanks

\end{document}
